\chapter[Introdução]{Introdução}
\addcontentsline{toc}{chapter}{Introdução}

O acesso à informação é um fator essencial para a boa experiência acadêmica dos estudantes universitários. Na Universidade de Brasília (UnB), informações institucionais relevantes, como prazos acadêmicos, procedimentos administrativos, normas para trancamento de disciplinas, emissão de documentos, entre outros, encontram-se frequentemente dispersas em documentos oficiais, editais e regulamentos que, por vezes, são numerosos e complexos.                                                                        

Além de todas as dificuldades já esperadas para uma jornada acadêmica (como a alta demanda de disciplinas, dificuldades profissionais e até problemas de relacionamentos com colegas e professores), a incerteza e demora em obter informações básicas, como processos necessários para o cumprimento de Estágios e Trabalhos de Conclusão de Curso (TCC), acaba elevando o nível de insatisfação dos estudantes. 

Em um estudo realizado por Claudino~\cite{claudino2023chatbot}, com relação a Faculdade de Ciência e Tecnologia em Engenharias (FCTE), chegou-se a conclusão que os estudantes da instituição possuem dificuldades relacionadas à utilização de portais acadêmicos, como excesso de etapas, estruturas desorganizadas e dificeis de entender e informações desatualizadas.

Além disso, conforme destacado por \cite{lorena2019informacao}, a qualidade das informações acadêmicas e administrativas disponibilizadas aos estudantes exerce impacto direto sobre a forma como estes organizam sua rotina ao longo do curso. Esse aspecto pode atuar tanto como fator de apoio quanto como obstáculo à jornada acadêmica, influenciando não apenas o desempenho individual, mas também a postura do discente em relação ao próprio processo de aprendizagem. Dessa maneira, a disponibilização clara e acessível de informações institucionais pode estimular uma mudança de atitude, conduzindo o estudante a assumir um papel mais ativo e protagonista em sua formação, em consonância com o apoio do corpo docente.

Ainda segundo Pan \cite{pan2024fragmentation}, a fragmentação de informações nas instituições de ensino superior gera um impacto negativo na formação, visto que um tempo relevante é gasto, por parte dos estudantes, para procurar as informações dispersas, prejudicando o entendimento e compreensão de diversos processos academicos. 

Tendo em mente que a fragmentação de informações espalhadas em várias plataformas é um problema para os estudantes, e isso também é uma realidade da FCTE, nota-se a chance de aprimorar a forma como os alunos encontram, acessam os dados e são atendidos pela instituição. Com o uso generalizado de aplicativos de comunicação Mobile, é válido a ideia de criar soluções tecnológicas que se encaixem em software já conhecidos, como WhatsApp, Telegram e outros aplicativos, adaptando-se na rotina dos discentes, e assim simplificando o acesso e suporte oferecido.

\section{Justificativa}

Diante dos estudos mencionados anteriormente, nota-se a oportunidade de buscar por alternativas tecnológicas que possam contribuir com a experiência acadêmica estudantil no que tange a suporte a informações institucionais administrativas da Universidade de Brasília.

Este trabalho, portanto, tem por finalidade o desenvolvimento de um sistema na forma de um bot que pode ser integrado a outras ferramentas, que terá como finalidade auxiliar os estudantes da FCTE no acesso simplificado a informações institucionais administrativas, assim como receber suporte administrativo do corpo da Universidade para solução de problemas acadêmicos. Para a construção do Bot, será utilizada uma base de dados construída a partir da extração e estruturação de documentos oficiais da UnB.

A proposta visa, portanto, melhorar a comunicação entre a instituição e seus estudantes, promovendo autonomia, agilidade e clareza no acesso às informações institucionais, e um local centralizado para os discentes receberem assistência estudantil. Para fundamentar a relevância da proposta, uma pesquisa foi realizada com estudantes da Universidade de Brasília do campus do Gama, cujos resultados contribuiram no desenvolvimento do software, principalmente na compreensão dos requisitos e espectativas da aplicação.

\section{Objetivos}

Esta seção tem como finalidade apresentar de forma clara e organizada os propósitos centrais deste trabalho de conclusão de curso. Os objetivos definidos servem como diretrizes para a condução da pesquisa e do desenvolvimento do sistema proposto, permitindo delimitar o escopo da solução e orientar as etapas metodológicas. Para isso, os objetivos estão divididos em objetivo geral, que expressa a meta principal do projeto, e objetivos específicos, que detalham as etapas e metas intermediárias necessárias para alcançar o resultado pretendido.

\subsection{Objetivo Geral}

Desenvolver uma ferramenta digital, na forma de um bot, com o objetivo de auxiliar estudantes da Universidade de Brasília (UnB) do campus do Gama (FCTE) no acesso rápido e facilitado a informações institucionais extraídas de documentos oficiais, assim como receber devido atendimento para problemas relacionados a sua tragetória acadêmica.

\subsection{Objetivos Específicos}

\begin{itemize}
    \item Identificar as principais dificuldades enfrentadas pelos estudantes da UnB no acesso a informações institucionais, por meio da aplicação de uma pesquisa exploratória;
    
    \item Levantar, organizar e estruturar conteúdos relevantes a partir de documentos oficiais da universidade, construindo uma base de dados de apoio ao bot;
    
    \item Projetar a arquitetura da ferramenta, definindo fluxos de conversa e mecanismos de busca e resposta;
    
    \item Implementar um bot funcional, com integração à base de dados construída;
    
    \item Disponibilizar as funcionalidades do Bot por meio de plataformas de comunicação (como WhatsApp, Telegram e outros);
    
    \item Validar a eficácia da ferramenta por meio de testes com usuários reais e coletar feedbacks para possíveis melhorias.
\end{itemize}

\section{Estrutura do Trabalho}

A seguir, apresenta-se a estrutura geral deste Trabalho de Conclusão de Curso:

\begin{enumerate}
    \item \textbf{Introdução}: Apresenta o tema, a delimitação do problema, os objetivos da pesquisa, a justificativa da proposta e a estrutura do documento.

    \item \textbf{Fundamentação Teórica}: Aborda os principais conceitos e estudos que embasam o projeto, incluindo tópicos como interação humano-chatbot, usabilidade, SLAs, ITIL v4, qualidade de software, design centrado no usuário e tecnologias aplicadas.

    \item \textbf{Metodologia}: Detalha os métodos utilizados na condução da pesquisa e desenvolvimento do projeto, incluindo a abordagem científica adotada, o processo de levantamento de dados, a aplicação do Design Thinking, o cronograma de atividades e os métodos de validação.

    \item \textbf{Resultados e Discussões}: Reúne e analisa os resultados obtidos durante o projeto. Esta seção está subdividida em:

    \begin{itemize}
        \item \textbf{Pesquisa com Estudantes da FCTE}: Apresenta os dados obtidos por meio da pesquisa aplicada, trazendo percepções sobre as principais dificuldades acadêmicas e as expectativas em relação ao uso de chatbots.
        
        \item \textbf{Benchmarking}: Analisa soluções similares adotadas por outras instituições, destacando funcionalidades, limitações e inspirações relevantes para o projeto.
        
        \item \textbf{Elicitação e Especificação de Requisitos}: Sistematiza os requisitos funcionais e não funcionais levantados a partir dos dados coletados, além de apresentar os critérios de aceitação e a priorização das funcionalidades.
        
        \item \textbf{Modelagem e Arquitetura da Solução}: Expõe a estrutura da solução proposta, incluindo a arquitetura da aplicação, os casos de uso e os principais componentes do sistema.
        
        \item \textbf{Prototipação}: Apresenta os protótipos de baixa fidelidade desenvolvidos como forma de ilustrar a interface e o fluxo de interação da solução com base nas funcionalidades planejadas.
    \end{itemize}
\end{enumerate}